\section{Perverse schobers on a disc}\label{section:adjaut}
We recall from \cite{kapranovPerverseSchobers2015} the following definition, which we will choose to refer to as abstract perverse schober on a disc.

\begin{defn}
    An abstract perverse schober on a disc with one singularity is a pair of stable categories, $\Phi,\Psi$ together with a spherical adjunction,
    \[S:\Phi\tofrom\Psi:R.\]
    An abstract perverse schober on a disc with $n$ singularities is a list of stable categories, $\Phi_1,\dots,\Phi_n,\Psi$, together with spherical adjunctions,
    \[S_i:\Phi_i\tofrom\Psi:R_i.\]
\end{defn}

\begin{notn}
    We denote by $\sfSph(n)$ the collection of abstract perverse schobers, which we later upgrade to a $\two$-category. We also write $\sfSph = \sfSph(1)$.
\end{notn}

We recall in \S\ref{subsection:enhancedFS} an operation defined in \cites{kapranovPerverseSchobersAlgebra2020,barbacoviCompositionTwoSpherical2021,harderPerverseSheavesCategories2019} which we will refer to as the enhanced Fukaya--Seidel functor,
\[\FS_{\sfSph}:\sfSph(n)\to\sfSph.\]
Geometrically, this operation can be thought of as merging the $n$ singularities into one. In the case that our abstract schober comes from a Lefschetz fibration, it is shown in \cite{christRelativeCalabiYauStructures2026}*{Theorem 6.1} that the vanishing cycle of $\FS_{\sfSph}$ agrees with the usual notion of Fukaya--Seidel category.

In \S\ref{subsection:simplicialadj} and \S\ref{subsection:simplicialsph}, we will upgrade $\sfSph(n)$ to a $\two$-category, $\FS_{\sfSph}$ to a functor of $\two$-categories, and the assignment
\[\sfSph(-):[n]\mapsto\sfSph(n)\]
to a simplicial $\two$-category in such a way that the unique active map $[1]\to [n]$ induces $\FS_{\sfSph}$. Moreover, the inert face maps $d_{n,0},d_{n,n}:\sfSph(n)\to\sfSph(n-1)$ are given by forgetting the first and last spherical adjunctions respectively so that $\sfSph(-)$ is evidently 1-Segal. In other words, we may view
\[\FS_{\sfSph}:\sfSph(2)\simeq\sfSph(1)\times_{\sfSph(0)}\sfSph(1)\to \sfSph(1)\]
as an associative binary operation.

In parallel, we construct a map of simplicial categories $\sfSph(-)\to\sfAut(-)$, and prove in Proposition~\ref{prop:activecartesian2} that over the active arrow $[1]\to[n]$, this induces a pullback square of $\two$-categories as follows.
\begin{equation}\label{equation:importantpullbacksquare}\begin{tikzcd}
    {\sfSph(n)} \pullback & \sfSph \\
    {\sfAut(n)} & \sfAut
    \arrow["{\FS_{\sfSph}}", from=1-1, to=1-2]
    \arrow[from=1-1, to=2-1]
    \arrow[from=1-2, to=2-2]
    \arrow["{\FS_{\sfAut}}", from=2-1, to=2-2]
\end{tikzcd}\end{equation}

In \S\ref{subsection:mutationsph}, we will use this to transfer the action of $\Br_n$ on $\sfAut(n)$ to one on $\sfSph(n)$, which is compatible with $\FS_{\sfSph}$, recovering a special case of the construction in the upcoming work \cite{acj}. We show that on objects, this agrees with the $\Br_n$-action on abstract perverse schobers defined in \cite{kapranovPerverseSchobers2015}*{(2.8)}. This leads to a definition of a $\two$-category of non-abstract perverse schobers, $\sftwoP(\bbC,R)$, together with a projection functor
\[\sftwoP(\bbC,R)\to\olsftwoP(\bbC,R)\]
and an enhanced Fukaya--Seidel functor
\[\FS_{\sftwoP}:\sftwoP(\bbC,R)\to \sftwoP(\bbC,0).\]

We record the following restatement of Proposition~\ref{prop:activecartesian2} for non-abstract schobers.
\begin{thm}\label{thm:merge}
    Let $R\subset\bbC$ be a finite subset. Then, there is a natural pullback square of $\two$-categories as follows.
    \[\begin{tikzcd}
        {\sftwoP(\bbC,R)} \pullback & \sftwoP(\bbC,0) \\
        {\olsftwoP(\bbC,R)} & \olsftwoP(\bbC,0)
        \arrow["{\FS_{\sftwoP}}", from=1-1, to=1-2]
        \arrow[from=1-1, to=2-1]
        \arrow[from=1-2, to=2-2]
        \arrow["{\FS_{\olsftwoP}}", from=2-1, to=2-2]
    \end{tikzcd}\]
\end{thm}

Finally, in \S\ref{subsection:fourier}, we will explain how Theorem~\ref{thm:merge} can be interpreted as the Fourier transform of perverse schobers studied in \cite{kapranovPerverseSchobersAlgebra2020}.

\begin{eg}\label{eg:waldhausen2}
    We continue with notation from Example~\ref{eg:waldhausen1}, where we constructed an object $(\frakS;S_{n-1}(F),T')$ of $\sfAut(n)$ from a spherical adjunction $F\adjto G$. Moreover, $T'$ was the twist of a spherical adjunction $F'\adjto G'$, so the square (\ref{equation:importantpullbacksquare}) implies that our object in fact lifts to $\sfSph(n)$.

    In \cite{gargGITRootStacks2026}, the construction $\sfSph\to \sfSph$ sending $F\adjto G\mapsto F'\adjto G'$ is interpreted as a pushforward
    \[f_*:\sftwoP(\bbC_z,0)\to \sftwoP(\bbC_y,0)\]
    along the map $f:\bbC_z\xto{z^n} \bbC_y$.

    We may perturb $z^n$ by a linear term so that it is Morse, $f_x(z) := z^n-xz$, where $x\in\bbC^\times$. This Morse function has $n$ critical values $R_x = \{y=0, y^{n-1} = x^n\}$. We interpret the lift to $\sfSph(n)$ as a factorization,
    \[\sftwoP(\bbC_z,0)\xto{(f_x)_*} \sftwoP(\bbC_y,R_x)\xto{\FS_{\sftwoP}} \sftwoP(\bbC_y,0).\]
    
    In Example~\ref{eg:waldhausen3}, we will explain how we may consider all $x$, including $x=0$ at once, and package this into a perverse schober on $\bbC^2_{x,y}$.
\end{eg}

\subsection{The enhanced Fukaya--Seidel functor}\label{subsection:enhancedFS}
In \cite{barbacoviCompositionTwoSpherical2021} and \cite{kapranovPerverseSchobersAlgebra2020}*{\S 3.5}, it is shown that given an object $(\Phi_i,\Psi,S_i,R_i)\in\sfSph(n)$, we may produce a category $\Phi$ called its Fukaya--Seidel category, equipped with an SOD of the form $\langle\Phi_1,\dots,\Phi_n\rangle$. This category is part of a spherical adjunction
\[S:\Phi\tofrom\Psi:R,\]
giving us a mapping of objects that we denote $\FS_{\sfSph}:\sfSph(n)\to\sfSph$, and that we refer to as the enhanced Fukaya--Seidel functor. Moreover, there are identifications $S\circ I_i\simeq S_i$.

All together, we have an operation,
\[\opname{Ind}:\begin{tikzcd}[ampersand replacement=\&]
	\& \Psi \& \\
	{\Phi_1} \& \cdots \& {\Phi_n}
	\arrow["{R_1}", shift left, from=1-2, to=2-1]
	\arrow["{R_n}", shift left, from=1-2, to=2-3]
	\arrow["{S_1}", shift left, from=2-1, to=1-2]
	\arrow["{S_n}", shift left, from=2-3, to=1-2]
\end{tikzcd}
\mapsto
% https://q.uiver.app/#q=WzAsNSxbMSwxLCJcXGxhbmdsZVxcUGhpXzEsXFxkb3RzLFxcUGhpX25cXHJhbmdsZSJdLFswLDIsIlxcUGhpXzEiXSxbMiwyLCJcXFBoaV9uIl0sWzEsMCwiXFxQc2kiXSxbMSwyLCJcXGNkb3RzIl0sWzAsMSwiSV8xXlIiLDIseyJvZmZzZXQiOjF9XSxbMSwwLCJJXzEiLDIseyJvZmZzZXQiOjF9XSxbMCwyLCJJX25eUiIsMix7Im9mZnNldCI6MX1dLFsyLDAsIklfbiIsMix7Im9mZnNldCI6MX1dLFszLDAsIlIiLDIseyJvZmZzZXQiOjF9XSxbMCwzLCJTIiwyLHsib2Zmc2V0IjoxfV1d
\begin{tikzcd}[ampersand replacement=\&]
	\& \Psi \& \\
	\& {\langle\Phi_1,\dots,\Phi_n\rangle} \\
	{\Phi_1} \& \cdots \& {\Phi_n}
	\arrow["R"', shift right, from=1-2, to=2-2]
	\arrow["S"', shift right, from=2-2, to=1-2]
	\arrow["{I_1^R}"', shift right, from=2-2, to=3-1]
	\arrow["{I_n^R}"', shift right, from=2-2, to=3-3]
	\arrow["{I_1}"', shift right, from=3-1, to=2-2]
	\arrow["{I_n}"', shift right, from=3-3, to=2-2]
\end{tikzcd}\]

We will now construct the operation $\opname{Ind}$ carefully for $n=2$. We will find it helpful to perform the construction without any sphericity/invertibility conditions so that we obtain a mapping of objects
\[\opname{Ind}:\sfAdj(2)\to \sfSOD(2)\times_{\sfSt_k}\sfAdj.\]

Suppose that $S_i:\Phi_i\tofrom\Psi:R_i$ for $i=1,2$ are two adjunctions, defining an object of $\sfAdj(2)$. Define $\Phi_{12}$ to be the category of diagrams of the form,
\[\begin{tikzcd}
    {S_2\phi_2} & \\
    {S_1R_1S_2\phi_2} & {S_1\phi_1}
    \arrow[from=2-1, to=1-1]
    \arrow[from=2-1, to=2-2]
\end{tikzcd}\]
together with a lift of the horizontal arrow to a diagram $R_1S_2\phi_2\to \phi_1$ in $\Phi_1$, and such that the vertical arrow is a counit for $S_1\adjto R_1$.

\begin{rmk}\label{rmk:laxcolimit}
    This diagram category is described in a somewhat unwieldy manner. Observe that the objects of this diagram category can also be written as a triple $(\phi_1,\phi_2,\alpha:R_1S_2\phi_2\to \phi_1)$ where $\phi_i\in\Phi_i$. Thus this is equivalent to the lax colimit of diagram $\Phi_2\xto{F_{12}}\Phi_1$ in the spirit of \cite{christLaxAdditivity2025}. In the language of \cite{kapranovPerverseSchobersAlgebra2020}, it is the category of modules for a unitriangular monad. That said, the diagram category allows for a nice description of $S_{12}$ below.
\end{rmk}

Consider the pair of maps
\[S_{12}:\Phi_{12}\tofrom\Psi:R_{12},\]
where 
\[S_{12}\left(\begin{tikzcd}
    {S_2\phi_2} & \\
    {S_1R_1S_2\phi_2} & {S_1\phi_1}
    \arrow[from=2-1, to=1-1]
    \arrow[from=2-1, to=2-2]
\end{tikzcd}\right) := \colim\left(\begin{tikzcd}
    {S_2\phi_2} & \\
    {S_1R_1S_2\phi_2} & {S_1\phi_1}
    \arrow[from=2-1, to=1-1]
    \arrow[from=2-1, to=2-2]
\end{tikzcd}\right),\]
and where
\[R_{12}(\psi) := \begin{tikzcd}
    {S_2R_2\psi} & \\
    {S_1R_1S_2R_2\psi} & {S_1R_1\psi}
    \arrow[from=2-1, to=1-1]
    \arrow[from=2-1, to=2-2]
\end{tikzcd}\]
where the horizontal arrow is the counit for $S_2\adjto R_2$.

\begin{lem}\label{lem:twocats}
    There is a natural adjunction, $S_{12}\adjto R_{12}$.
\end{lem}

\begin{proof}
    Given a diagram $\phi_{12}$ as above, we have functorial identifications
    \begin{align*}
        \Hom(S_{12}\phi_{12},\psi)&\simeq \Hom(S_1\phi_1,\psi)\times_{\Hom(S_1R_1S_2\phi_2,\psi)} \Hom(S_2\phi_2,\psi)\\
        &\simeq \Hom(\phi_1,R_1\psi)\times_{\Hom(R_1S_2\phi_2,R_1\psi)} \Hom(\phi_2,R_2\psi)\\
        &\simeq \Hom(\phi_{12},R_{12}\psi).
    \end{align*}
\end{proof}

\begin{lem}\label{lem:twocatssod}
    The category $\Phi_{12}$ admits a right-admissible SOD via the following subcategories. Here, we denote elements of $\Phi_{12}$ by triples $(\phi_1,\phi_2,\alpha:R_1S_2\phi_2\to \phi_1)$, as discussed in Remark~\ref{rmk:laxcolimit}.
    \begin{align*}
        I_1: \Phi_1&\to \Psi\\
        \phi_1&\mapsto (\phi_1,0,0)\\
        I_2: \Phi_2&\to \Psi\\
        \phi_2&\mapsto (R_1S_2\phi_2,\phi_2,\id_{R_1S_2\phi_2})
    \end{align*}
    The right-adjoints are given by projections $I_i^R:(\phi_1,\phi_2,\alpha)\mapsto \phi_i$. Moreover, there are canonical identifications $S_{12}\circ I_i\simeq S_i$.
\end{lem}

\begin{proof}
    Explicit calculation.
\end{proof}

Recall that for $\calF\in\sfAdj(2)$, we may recover $\calF$ from $\opname{Ind}(\calF)$. Thus $\opname{Ind}$ is an injective mapping. The following computes the image of this mapping.

\begin{lem}\label{lem:twocatsequivalence}
    The image of $\opname{Ind}:\sfAdj(2)\to\sfSOD(2)\times_{\sfSt_k}\sfAdj$ consists of adjunctions $S:\Phi\tofrom\Psi:R$ equipped with a 2-term SOD $\frakS$ of $\Phi$ such that $\coCone(1\to RS)$ is weakly conjugate-compatible with respect to $\frakS$.
\end{lem}

\begin{proof}
    Given an object $\calF\in\sfAdj(2)$, we show that $\coCone(\id\to R_{12}S_{12})$ is weakly conjugate-compatible with respect to the SOD $\frakS$ on $\Phi_{12}$ as constructed in Lemma~\ref{lem:twocatssod}. This is by direct calculation - see that
    \[I_1^R\coCone(\id\to R_{12}S_{12})I_2\simeq \coCone(I_1^RI_2\to R_1S_2) = 0.\]

    We now prove the other containment. We consider an object of $\sfSOD(2)\times_{\sfSt_k}\sfAdj$ given as follows.
    \[\begin{tikzcd}[ampersand replacement=\&]
        \& \Psi \& \\
        \& {\langle\Phi_1,\Phi_2\rangle} \\
        {\Phi_1} \&  \& {\Phi_n}
        \arrow["R"', shift right, from=1-2, to=2-2]
        \arrow["S"', shift right, from=2-2, to=1-2]
        \arrow["{I_1^R}"', shift right, from=2-2, to=3-1]
        \arrow["{I_2^R}"', shift right, from=2-2, to=3-3]
        \arrow["{I_1}"', shift right, from=3-1, to=2-2]
        \arrow["{I_2}"', shift right, from=3-3, to=2-2]
    \end{tikzcd}\]
    Consider the object $\calF\in\sfAdj(2)$ given by $(\Phi_i,\Psi,SI_i,I_i^RR)$. Because $\coCone(\id\to RS)$ is conjugate-compatible with respect to $\frakS$, the natural map
    \[I_1^RI_2\to I_1^RRSI_2\]
    is an equivalence. Using this, we deduce that $\Phi_{12}$ is equivalent to the category of diagrams of form
    \[\begin{tikzcd}
        {SI_2\phi_2} & \\
        {SI_1I_1^RI_2\phi_2} & {SI_1\phi_1}
        \arrow[from=2-1, to=1-1]
        \arrow[from=2-1, to=2-2]
    \end{tikzcd}\]
    together with a lift of the horizontal arrow to a diagram $I_1I_1^RI_2\phi_2\to I_1\phi_1$ in $\Phi$, and such that the vertical arrow is a counit for $I_1\adjto I_1^R$. Thus $S\circ-$ can be factored out and this is equivalent to the diagram category described in Lemma~\ref{lem:twosubcats}, so that $\Phi_{12}\simeq \Phi$. Moreover, under this identification, the adjunction of Lemma~\ref{lem:twocats} is identified with $S\adjto R$ and the SOD of Lemma~\ref{lem:twocatssod} is identified with $\frakS$.
\end{proof}

\begin{lem}\label{lem:twocatsspherical}
    The image of $\opname{Ind}$ restricted to $\sfSph(2)$ consists of \textit{spherical} adjunctions $S:\Phi\tofrom\Psi:R$ equipped with a 2-term SOD $\frakS$ of $\Phi$ such that $\coCone(\id\to RS)$ is conjugate-compatible with respect to $\frakS$.
\end{lem}

\begin{proof}
    We first prove that given $\calF\in\sfSph(2)$, $\opname{Ind}(\calF)$ has the claimed properties. The case $n=2$ of \cite{kapranovPerverseSchobersAlgebra2020}*{Proposition 3.5.6} shows that $S_{12}\adjto R_{12}$ is spherical. This also follows by a direct calculation. See that 
    \[I_i^R\coCone(\id\to R_{12}S_{12})I_i\simeq \coCone(\id\to R_iS_i)\]
    is invertible. Together with Lemma~\ref{lem:twocatsequivalence}, this proves conjugate-compatibility.

    For the other containment, we must show that if $(\frakS;\Phi,\Psi,S,R)\in \sfSOD(2)\times_{\sfSt_k}\sfAdj$ has the claimed properties, then $SI_i:\Phi_i\tofrom \Psi:R_i$ is spherical. The computation above shows that $\coCone(\id\to R_iS_i)$ is invertible for each $i$. We extract from the proof of \cite{kapranovPerverseSchobersAlgebra2020}*{Proposition 3.5.6}, or show by direct calculation, that
    \[\Cone(SR\to \id)\simeq \Cone(S_1R_1\to \id)\Cone(S_2R_2\to \id).\]
    Since $\Cone(SR\to \id)$ is invertible, this shows that $\Cone(S_1R_1\to \id)$ has a right-inverse, and $\Cone(S_2R_2\to \id)$ has a left-inverse.

    By Theorem~\ref{thm:inftyadmissible}, we may mutate $\frakS$ to $\sigma\cdot\frakS = (\Phi_2',\Phi_1)$, and we denote the new inclusion by $I_2'$. This defines a new object $(\sigma\cdot\frakS;\Phi,\Psi,S,R)$ and the same argument above shows that $\Cone(S_1R_1\to \id)$ has a right-inverse. So we conclude that $\Cone(S_1R_1\to \id)$ is invertible. The other mutation shows that $\Cone(S_2R_2\to \id)$ is invertible.
\end{proof}

\begin{rmk}
    A version of this (in one direction) appears in \cite{halpern-leistnerAutoequivalencesDerivedCategories2016}*{Theorem 4.4}. This statement contains a condition equivalent to our conjugate-compatibility condition.
\end{rmk}

For $n>2$, we will construct $\opname{Ind}$ by abstract means in Theorem~\ref{prop:bdjisadj}.

\subsection{\texorpdfstring{$\sfAdj(n)$}{Adj(n)} as a simplicial \texorpdfstring{$\two$}{2}-category}\label{subsection:simplicialadj}
We construct the simplicial $\two$-category $\sfAdj(-)$. We do so by constructing an a priori different simplicial $\two$-category $\sfBdj(-)$, and then comparing $\sfBdj(n)$ to $\sfAdj(n)$.

Consider the functor
\[\sfAdj \to \sfEnd,\]
given by $(\Phi,\Psi,S,R)\mapsto (\Phi,\coCone(\id\to RS))$. Consider the simplicial $\two$-category
\[\sfBdj'(-) := \sfEnd'(-)\times_{\sfEnd}\sfAdj.\]
The objects of $\sfBdj'(n)$ are given by a category $\Phi$ with a partial SOD $\frakS$ and an adjunction $(\Phi,\Psi,S,R)$ such that $\coCone(\id\to RS)$ is weakly conjugate-compatible with respect to $\frakS$.

For each $n$, consider the full sub-$\two$-category $\opname{I}_n:\sfBdj(n)\inclto \sfBdj'(n)$ given by those objects where the partial SOD is an SOD.

\begin{lem}\label{lem:rightadjointbdj}
    The inclusion functor $\opname{I}_n:\sfBdj(n) \inclto \sfBdj'(n)$ admits a right adjoint $\opname{I}_n^R$, which is given on objects by
    \[(\frakS;\Phi,\Psi,S,R)\mapsto (\frakS;\langle\Phi_1,\dots,\Phi_n\rangle,\Psi,SI,I^RR),\]
    where $I:\langle\Phi_1,\dots,\Phi_n\rangle\inclto \Phi$.
\end{lem}

\begin{proof}
    This is proven similarly to Lemma~\ref{lem:keyrightadjoint} and Lemma~\ref{lem:rightadjointend}, by reducing to the case $n=1$.
\end{proof}

\begin{cor}\label{prop:simplicialbdj}
    The assignment $[n]\mapsto \sfBdj(n)$ upgrades to a simplicial $\two$-category. For $0\leq i\leq n$, the degeneracy map $s_{n,i}$ inserts a $0$ between $\Phi_i$ and $\Phi_{i+1}$;
    \begin{align*}
        s_{n,i}:\sfBdj(n)&\to\sfBdj(n+1)\\
        (\frakS;\Phi,\Psi,S,R)&\mapsto ((\Phi_1,\dots,\Phi_i,0,\Phi_{i+1},\dots\Phi_n);\Phi,\Psi,S,R).
    \end{align*}
    The inert face map $d_{n,0}$ removes $\Phi_1$ from the SOD to get $\Phi_{2,\dots,n} := \langle\Phi_2,\dots,\Phi_n\rangle\xto{I_{2,\dots,n}}\Phi$;
    \begin{align*}
        d_{n,0}:\sfBdj(n)&\to\sfBdj(n-1)\\
        (\frakS;\Phi,\Psi,S,R)&\mapsto ((\Phi_2,\dots,\Phi_n);\Phi_{2,\dots,n},\Psi,SI_{2,\dots,n},I_{2,\dots,n}^RR),
    \end{align*}
    and the other inert face map $d_{n,n}$ similarly removes $\Phi_n$ from the SOD. For $0<i<n$, the active face map is given by merging $\Phi_{i}$ and $\Phi_{i+1}$ in the SOD into $\Phi_{i,i+1} := \langle\Phi_i,\Phi_{i+1}\rangle$;
    \begin{align*}
        d_{n,i}:\sfBdj(n)&\to\sfBdj(n-1)\\
        (\frakS;\Phi,\Psi,S,R)&\mapsto ((\Phi_1,\dots,\Phi_{i-1},\Phi_{i,i+1},\Phi_{i+2},\dots\Phi_n);\Phi,\Psi,S,R).
    \end{align*}
\end{cor}

\begin{proof}
    We apply Lemma~\ref{lem:descendingsimplicialobjects}. We omit the verification that the resulting lax-commuting squares are commuting squares.
\end{proof}

\begin{prop}
    The simplicial $\two$-category $\sfBdj(-)$ is 2-Segal.
\end{prop}

\begin{proof}
    The strategy used in the proof of Proposition~\ref{prop:segalsod} works, using that the map $\sfBdj'(-)\to\sfBdj(-)$ is Cartesian over $\Delta_{\act}$.
\end{proof}

We now construct a projection $\sfBdj(-)\to\sfEnd(-)$ of simplicial $\two$-categories.

\begin{prop}\label{prop:simplicialbdjtoend}
    There exists a map $\pi(-): \sfBdj(-)\to\sfEnd(-)$ of simplicial $\two$-categories, which for each $n$ is given object-wise by
    \[(\frakS;\Phi,\Psi,R,S)\mapsto (\frakS;\Phi,\coCone(\id\to RS)).\]
\end{prop}

\begin{proof}
    We begin with the simplicial map $\pi'(-):\sfBdj'(-)\to\sfEnd'(-)$. We apply Lemma~\ref{lem:descendingsimplicialmaps} to define functors
    \[\pi(n):\sfBdj(n)\to \sfEnd(n)\]
    which are given on objects by the formulae stated. For each $n$, the relevant lax-commuting square evidently commutes so we obtain the desired map $\pi(-)$.
\end{proof}

\begin{prop}\label{prop:activecartesian4}
    The map $\pi(-): \sfBdj(-)\to\sfEnd(-)$ of simplicial $\two$-categories is Cartesian over $\Delta_{\act}$.
\end{prop}

\begin{proof}
    We are to show that for each $n$, the following commutative square is a pullback.
    \[\begin{tikzcd}[ampersand replacement=\&]
        {\sfBdj(n)} \& {\sfBdj(1)} \\
        {\sfEnd(n)} \& {\sfEnd(1)}
        \arrow[from=1-1, to=1-2]
        \arrow[from=1-1, to=2-1]
        \arrow[from=1-2, to=2-2]
        \arrow[from=2-1, to=2-2]
    \end{tikzcd}\]
    This is clear by unravelling the definitions.
\end{proof}

We now study $\sfAdj(n)$. Observe first that there is a natural identification,
\[\sfAdj(n)\simeq \sfAdj(1)\times_{\sfAdj(0)}\cdots\times_{\sfAdj(0)}\sfAdj(1),\]
where $\sfAdj(1)\to \sfAdj(0)$ is given by the functor $(\Phi,\Psi,S,R)\mapsto \Psi$. Thus $[n]\mapsto \sfAdj(n)$ organize into a presheaf of $\two$-categories over $\Delta_{\inert}$.

We will construct an equivalence, $\sfBdj(-)\resto{\Delta_{\inert}}\simeq\sfAdj(-)$, thereby upgrading $\sfAdj(-)$ to a 1-Segal simplicial $\two$-category.

For each fixed $n$, consider the map of objects $\sfBdj(n)\to\sfAdj(n)$ given by
\[(\frakS;\Phi,\Psi,S,R)\mapsto (\Phi_i,\Psi,SI_i,I_i^RR),\]
where $I_i:\Phi_I\inclto \Phi$ are the terms of the SOD. To construct this as a functor, we pass to diagram categories as follows.

Let $\frakn(n)$ be the ordinary 1-category with objects $1,\dots,n,\infty$ and arrows $i\to \infty$. Let $\frakh(n)$ be the ordinary 1-category with objects $1,\dots,n,0,\infty$ and arrows $i\to \infty$, $\infty\to 0$. There is a natural functor $\frakn(n)\to\frakh(n)$, depicted below.
% https://q.uiver.app/#q=WzAsMTEsWzAsMiwiMSJdLFsxLDIsIlxcY2RvdHMiXSxbMiwyLCJuIl0sWzEsMCwiXFxpbmZ0eSJdLFsyLDFdLFszLDFdLFszLDIsIjEiXSxbNCwxLCIwIl0sWzQsMiwiXFxjZG90cyJdLFs1LDIsIm4iXSxbNCwwLCJcXGluZnR5Il0sWzAsM10sWzIsM10sWzEsM10sWzQsNV0sWzYsN10sWzgsN10sWzksN10sWzcsMTBdXQ==
\[\begin{tikzcd}[ampersand replacement=\&]
	\& \infty \&\&\& \infty \& \\
	\&\& {} \& {} \& 0 \\
	1 \& \cdots \& n \& 1 \& \cdots \& n
	\arrow[from=2-3, to=2-4]
	\arrow[from=2-5, to=1-5]
	\arrow[from=3-1, to=1-2]
	\arrow[from=3-2, to=1-2]
	\arrow[from=3-3, to=1-2]
	\arrow[from=3-4, to=2-5]
	\arrow[from=3-5, to=2-5]
	\arrow[from=3-6, to=2-5]
\end{tikzcd}\]

By construction, $\sfAdj(n)\to \sfFun(\frakn(n),\sfSt_k)$ and likewise, $\sfBdj(n)\to \sfFun(\frakh(n),\sfSt_k)$ are locally full sub-$\two$-categories. Pullback along $\frakn(n)\to \frakh(n)$ induces the claimed functors $\sfBdj(n)\to\sfAdj(n)$.

Observe further that $\sfFun(\frakn(-),\sfSt_k)$, $\sfFun(\frakh(-),\sfSt_k)$ organize into presheaves of $\two$-categories over $\Delta_{\inert}$ because $\frakn(-)$ and $\frakh(-)$ organize into precosheaves. This induces the aforementioned presheaf $\sfAdj(-)$. It also induces the presheaf $n\mapsto \sfBdj(n)$, which agrees with $\sfBdj\resto{\Delta_{\inert}}$, and thus we have the desired map,
\[\opname{C}(-):\sfBdj(-)\resto{\Delta_{\inert}}\to\sfAdj(-).\]

\begin{prop}\label{prop:bdjisadj}
    The map $\opname{C}(n):\sfBdj(n)\to\sfAdj(n)$ is an equivalence for each $n$. Thus $\sfAdj(-)$ extends to a 1-Segal simplicial $\two$-category.
\end{prop}

\begin{proof}
    The case $n=0$ and $n=1$ are trivial. We prove the result by induction on $n$, reducing to the base case $n=2$, which for the moment we assume.
    
    Let $n\geq3$ and assume the result for all smaller cases. Consider the commuting diagram in $\Delta$ of the form
    % https://q.uiver.app/#q=WzAsNixbMCwwLCJbbl0iXSxbMCwxLCJbMl0iXSxbMSwxLCJbMV0iXSxbMSwwLCJbbi0xXSJdLFsyLDAsIltuLTJdIl0sWzIsMSwiWzBdIl0sWzIsMSwiYSJdLFs1LDJdLFs1LDRdLFs0LDNdLFsyLDMsImkiXSxbMSwwXSxbMywwXV0=
    \begin{equation}\label{equation:deltadiagram}\begin{tikzcd}[ampersand replacement=\&]
        {[n]}\pullback \& {[n-1]}\pullback \& {[n-2]} \\
        {[2]} \& {[1]} \& {[0]}
        \arrow[from=1-2, to=1-1]
        \arrow["j", from=1-3, to=1-2]
        \arrow[from=2-1, to=1-1]
        \arrow["i", from=2-2, to=1-2]
        \arrow["a", from=2-2, to=2-1]
        \arrow[from=2-3, to=1-3]
        \arrow[from=2-3, to=2-2]
    \end{tikzcd}
    \end{equation}
    where both squares are pushouts, $a$ is active, $i$ is the inert map with image $\{0,1\}$, and $j$ is the inert map with image $\{1,\dots,n-1\}$. Note that the outer square is also a pushout square involving only inert maps.

    By applying $\sfAdj(-)$ to the inert arrows, and $\sfBdj(-)$ to everything, this induces the following commutative diagram.
    % https://q.uiver.app/#q=WzAsMTIsWzEsMywiXFxzZkFkaigyKSJdLFsyLDMsIlxcc2ZBZGooMSkiXSxbMSwyLCJcXHNmQWRqKG4pIl0sWzIsMiwiXFxzZkFkaihuLTEpIl0sWzMsMiwiXFxzZkFkaihuLTIpIl0sWzMsMywiXFxzZkFkaigwKSJdLFswLDEsIlxcc2ZCZGooMikiXSxbMSwxLCJcXHNmQmRqKDEpIl0sWzAsMCwiXFxzZkJkaihuKSJdLFsxLDAsIlxcc2ZCZGoobi0xKSJdLFsyLDAsIlxcc2ZCZGoobi0yKSJdLFsyLDEsIlxcc2ZCZGooMCkiXSxbMiwwXSxbMywxXSxbMyw0XSxbNCw1XSxbMSw1XSxbMiwzLCJcXG9wbmFtZXtZfSIsMix7InN0eWxlIjp7ImJvZHkiOnsibmFtZSI6ImRhc2hlZCJ9fX1dLFswLDEsIlxcb3BuYW1le1h9IiwyLHsic3R5bGUiOnsiYm9keSI6eyJuYW1lIjoiZGFzaGVkIn19fV0sWzYsN10sWzgsNl0sWzgsOV0sWzksMTBdLFsxMCwxMV0sWzcsMTFdLFs5LDddLFs2LDBdLFs4LDJdLFs3LDFdLFs5LDNdLFsxMSw1XSxbMTAsNF1d
    \[\begin{tikzcd}[ampersand replacement=\&]
        {\sfBdj(n)}\pullback \& {\sfBdj(n-1)} \& {\sfBdj(n-2)} \& \\
        {\sfBdj(2)} \& {\sfBdj(1)} \& {\sfBdj(0)} \\
        \& {\sfAdj(n)} \& {\sfAdj(n-1)}\pullback \& {\sfAdj(n-2)} \\
        \& {\sfAdj(2)} \& {\sfAdj(1)} \& {\sfAdj(0)}
        \arrow[from=1-1, to=1-2]
        \arrow[from=1-1, to=2-1]
        \arrow[from=1-1, to=3-2]
        \arrow[from=1-2, to=1-3]
        \arrow[from=1-2, to=2-2]
        \arrow[from=1-2, to=3-3]
        \arrow[from=1-3, to=2-3]
        \arrow[from=1-3, to=3-4]
        \arrow[from=2-1, to=2-2]
        \arrow[from=2-1, to=4-2]
        \arrow[from=2-2, to=2-3]
        \arrow[from=2-2, to=4-3]
        \arrow[from=2-3, to=4-4]
        \arrow["{\opname{Y}}"', dashed, from=3-2, to=3-3]
        \arrow[from=3-2, to=4-2]
        \arrow[from=3-3, to=3-4]
        \arrow[from=3-3, to=4-3]
        \arrow[from=3-4, to=4-4]
        \arrow["{\opname{X}}"', dashed, from=4-2, to=4-3]
        \arrow[from=4-3, to=4-4]
    \end{tikzcd}\]
    
    To navigate this diagram, we will interpret the diagonal arrows as coming out of the page. The back-left square is a pullback by the 2-Segal property of $\sfBdj(-)$, and front-right square is a pullback by the 1-Segal property of $\sfAdj(-)$. The diagonal maps come from $\opname{C}(-)$ and all but $\opname{C}(n)$ is an equivalence by inductive hypothesis. These equivalences induce the dotted map denoted $\opname{X}$. Finally, the dotted map $\opname{Y}$ can be completed using the pullback square on the front-right.

    The front-outer square is a pullback by the 1-Segal property of $\sfAdj(-)$ applied to the outer square of (\ref{equation:deltadiagram}). Thus the front-left square is a pullback, and we deduce that $\opname{C}(n)$ is an equivalence as desired.

    It remains to prove the base case $n=2$. We embed our two $\two$-categories into slightly larger $\two$-categories as follows. Let $\frakn'$ and $\frakh'$ be the small $\two$-categories as shown, together with a map $f:\frakn'\to\frakh'$.
    % https://q.uiver.app/#q=WzAsOSxbMywyLCIxIl0sWzQsMCwiXFxpbmZ0eSJdLFs1LDIsIjIiXSxbMCwyLCIxIl0sWzIsMiwiMiJdLFsxLDAsIlxcaW5mdHkiXSxbMiwxXSxbMywxXSxbNCwxLCIwIl0sWzIsMF0sWzMsNV0sWzQsNV0sWzQsM10sWzYsNywiZiJdLFswLDhdLFsyLDhdLFs4LDFdLFszLDExLCIiLDAseyJzaG9ydGVuIjp7InRhcmdldCI6MjB9fV0sWzAsMTUsIiIsMCx7InNob3J0ZW4iOnsidGFyZ2V0IjoyMH19XV0=
    \[\begin{tikzcd}[ampersand replacement=\&]
        \& \infty \&\&\& \infty \& \\
        \&\& {} \& {} \& 0 \\
        1 \&\& 2 \& 1 \&\& 2
        \arrow["f", from=2-3, to=2-4]
        \arrow[from=2-5, to=1-5]
        \arrow[from=3-1, to=1-2]
        \arrow[""{name=0, anchor=center, inner sep=0}, from=3-3, to=1-2]
        \arrow[from=3-3, to=3-1]
        \arrow[from=3-4, to=2-5]
        \arrow[""{name=1, anchor=center, inner sep=0}, from=3-6, to=2-5]
        \arrow[from=3-6, to=3-4]
        \arrow[between={0}{0.8}, Rightarrow, from=3-1, to=0]
        \arrow[between={0}{0.8}, Rightarrow, from=3-4, to=1]
    \end{tikzcd}\]
    We may embed $\opname{J}:\sfAdj(2)\to \sfFun(\frakn',\sfSt_k)$ as a locally full sub-$\two$-category.
    % https://q.uiver.app/#q=WzAsOCxbMywyLCJcXFBoaV8xIl0sWzUsMiwiXFxQaGlfMiJdLFs0LDAsIlxcUHNpIl0sWzAsMiwiXFxQaGlfMSJdLFsxLDAsIlxcUHNpIl0sWzIsMiwiXFxQaGlfMiJdLFszLDFdLFsyLDFdLFswLDIsIlNfMSJdLFsxLDIsIlNfMiIsMl0sWzEsMCwiUl8xU18yIl0sWzMsNCwiU18xIiwwLHsib2Zmc2V0IjotMX1dLFs1LDQsIlNfMiIsMCx7Im9mZnNldCI6LTF9XSxbNCwzLCJSXzEiLDAseyJvZmZzZXQiOi0xfV0sWzQsNSwiUl8yIiwwLHsib2Zmc2V0IjotMX1dLFs3LDYsIiIsMCx7InN0eWxlIjp7InRhaWwiOnsibmFtZSI6Im1hcHMgdG8ifX19XSxbMCw5LCIiLDAseyJzaG9ydGVuIjp7InRhcmdldCI6MjB9fV1d
    \[\begin{tikzcd}[ampersand replacement=\&]
        \& \Psi \&\&\& \Psi \& \\
        \&\& {} \& {} \\
        {\Phi_1} \&\& {\Phi_2} \& {\Phi_1} \&\& {\Phi_2}
        \arrow["{R_1}", shift left, from=1-2, to=3-1]
        \arrow["{R_2}", shift left, from=1-2, to=3-3]
        \arrow[maps to, from=2-3, to=2-4]
        \arrow["{S_1}", shift left, from=3-1, to=1-2]
        \arrow["{S_2}", shift left, from=3-3, to=1-2]
        \arrow["{S_1}", from=3-4, to=1-5]
        \arrow[""{name=0, anchor=center, inner sep=0}, "{S_2}"', from=3-6, to=1-5]
        \arrow["{R_1S_2}", from=3-6, to=3-4]
        \arrow[between={0}{0.8}, Rightarrow, from=3-4, to=0]
    \end{tikzcd}\]
    Indeed, it is the sub-$\two$-category on objects $i\mapsto \calA_i$ such that the arrows $F_i:\calA_i\to \calA_\infty$ are left adjoints, and such that the lax-commuting triangle becomes a commuting triangle once we pass to the right adjoint of $F_1$. Moreover, given two objects $i\mapsto \calA_i,\calB_i$, we must restrict to those morphisms which are adjointable with over the arrows $i\to \infty$ (see \S\ref{subsection:mndadj}). Likewise, we may embed $\opname{K}:\sfBdj(2)\to \sfFun(\frakh',\sfSt_k)$ as a locally full sub-$\two$-category.

    Observe now that the pullback functor $f^*:\sfFun(\frakh',\sfSt_k)\to \sfFun(\frakn',\sfSt_k)$ restricts to these locally full sub-$\two$-categories. Indeed, given an object $\calF' := (\Phi_1',\Phi_2';\Phi',\Psi',S',R')$ with embeddings $I_i':\Phi_i'\inclto \Phi$, the image under $f^*\circ\opname{K}$ is as follows.
    % https://q.uiver.app/#q=WzAsMyxbMCwyLCJcXFBoaV8xJyJdLFsyLDIsIlxcUGhpXzInIl0sWzEsMCwiXFxQc2knIl0sWzAsMiwiU0lfMSciXSxbMSwyLCJTSV8yJyIsMl0sWzEsMCwiSV8xJ15SSV8yJyJdLFswLDQsIiIsMCx7InNob3J0ZW4iOnsidGFyZ2V0IjoyMH19XV0=
    \[\begin{tikzcd}[ampersand replacement=\&]
        \& {\Psi'} \& \\
        \\
        {\Phi_1'} \&\& {\Phi_2'}
        \arrow["{SI_1'}", from=3-1, to=1-2]
        \arrow[""{name=0, anchor=center, inner sep=0}, "{SI_2'}"', from=3-3, to=1-2]
        \arrow["{I_1'^RI_2'}", from=3-3, to=3-1]
        \arrow[between={0}{0.8}, Rightarrow, from=3-1, to=0]
    \end{tikzcd}\]
    For this to factor through $\opname{J}$ we require that $I_1'^RRSI_2'\to I_1'^RI_2'$ is an equivalence, which follows as its cocone is $I_1'^R\coCone(1\to R'S')I_2'=0$. The verification that morphisms of objects also factor through $\sfAdj(2)$ is immediate.
    
    The resulting functor is identified with $\opname{C}(2)$ (by considering the functors $\frakn\to\frakn',\frakh\to\frakh'$), thus we have the following diagram of $\two$-categories.
    \[\begin{tikzcd}[ampersand replacement=\&]
        {\sfBdj(2)} \& {\sfAdj(2)} \\
        {\sfFun(\frakh',\sfSt_k)} \& {\sfFun(\frakn',\sfSt_k)}
        \arrow["\opname{C}(2)", from=1-1, to=1-2]
        \arrow["\opname{K}", from=1-1, to=2-1]
        \arrow["\opname{J}", from=1-2, to=2-2]
        \arrow["f^*", from=2-1, to=2-2]
    \end{tikzcd}\]
    Importantly, the vertical maps are locally fully faithful - i.e. 
    By Theorem~\ref{thm:laxfibrations}, $f^*$ admits a left adjoint $f_!$, which moreover admits an explicit description. Given $\calF\in\sfAdj(2)$, we can compute $f_! \circ \opname{J}(\calF)$ and see that the object agrees with $\opname{K}(\opname{Ind}(\calF))$, where $\opname{Ind}(\calF)$ was constructed on the level of objects in \S\ref{subsection:enhancedFS}. Similarly, we see that given a map $P:\calF_1\to\calF_2$, the map $f_! \circ \opname{J}(P)$ also lifts along $\opname{K}$.

    Thus $\opname{Ind}$ canonically extends to a functor such that the following commutes.
    \[\begin{tikzcd}[ampersand replacement=\&]
        {\sfBdj(2)} \& {\sfAdj(2)} \\
        {\sfFun(\frakh',\sfSt_k)} \& {\sfFun(\frakn',\sfSt_k)}
        \arrow["{\opname{K}}"', from=1-1, to=2-1]
        \arrow["{\opname{Ind}}"', from=1-2, to=1-1]
        \arrow["{\opname{J}}", from=1-2, to=2-2]
        \arrow["{f_!}", from=2-2, to=2-1]
    \end{tikzcd}\]

    For notational convenience, we denote the four $\two$-categories in the diagram above by $\sfA,\sfB,\sfN',\sfH'$. Let $\calF\in\sfA=\sfAdj(2)$, $\calF'\in\sfB=\sfBdj(2)$. Then, we have a diagram of mapping categories as follows.
    \[\begin{tikzcd}[ampersand replacement=\&]
        {\sfB(\calF',\opname{Ind}(\calF))} \& {\sfA(\opname{C}(2)(\calF'),\calF)} \\
        {\sfH'(\opname{K}(\calF'),f_!\opname{J}(\calF))} \& {\sfN'(f^*\opname{K}(\calF'),\opname{J}(\calF))}
        \arrow["{j'}"', hook, from=1-1, to=2-1]
        \arrow["j", hook, from=1-2, to=2-2]
        \arrow["\sim", tail reversed, from=2-2, to=2-1]
    \end{tikzcd}\]
    The vertical maps are full subcategories because $\opname{J},\opname{K}$ are locally fully faithful. Verifying that the subcategories agree under the adjunction $f_!\adjto f^*$, we obtain an adjunction $\opname{Ind}\adjto\opname{C}(2)$. Moreover, $\opname{Ind}$ is 2-fully faithful because $f_!$ is.

    It remains to prove that $\opname{Ind}$ is essentially surjective. This is the calculation outlined in Lemma~\ref{lem:twocatsequivalence}.
\end{proof}

\subsection{\texorpdfstring{$\sfSph(n)$}{Sph(n)} as a simplicial \texorpdfstring{$\two$}{2}-category}\label{subsection:simplicialsph}

We now reinstate the sphericity conditions and construct the simplicial $\two$-category $\sfSph(-)$. We then construct an action of the braid group on $\sfSph(n)$ and show that this agrees with the braid group action defined at the level of objects in \cite{kapranovPerverseSchobers2015}. We use this to define the $\two$-category $\sftwoP(\bbC,R)$ of perverse schobers on $(\bbC,R)$.

\begin{prop}\label{prop:segalsph}
    The sub-$\two$-categories $\sfSph(n)\inclto \sfAdj(n)$ define a sub-simplicial $\two$-category which is 1-Segal.
\end{prop}

\begin{proof}
    It follows from the calculation of \cite{barbacoviCompositionTwoSpherical2021} or \cite{kapranovPerverseSchobersAlgebra2020}*{Proposition 3.5.6} that the ``monoidal product''
    \[d_{2,1}:\sfAdj(2)\to\sfAdj\]
    sends abstract perverse schobers to abstract perverse schobers. The same is clearly true for the ``monoidal unit''
    \[d_{0,0}:\sfAdj(0)\to\sfAdj.\]
    Moreover, there is an identification,
    \[\sfSph(n)\simeq \sfSph(1)\times_{\sfSph}\cdots\times_{\sfSph}\sfSph(1)\]
    compatibly with the analogous identification for $\sfAdj(n)$. The result now follows.
\end{proof}

\begin{prop}\label{prop:activecartesian2}
    There is a natural map of simplicial $\two$-categories $\pi(-):\sfSph(-)\to\sfAut(-)$ which is Cartesian over $\Delta_{\act}$. In particular, for each $n$, there is a natural pullback square as follows.
    \[\begin{tikzcd}
        {\sfSph(n)} \pullback & \sfSph \\
        {\sfAut(n)} & \sfAut
        \arrow["{\FS_{\sfSph}}", from=1-1, to=1-2]
        \arrow[from=1-1, to=2-1]
        \arrow[from=1-2, to=2-2]
        \arrow["{\FS_{\sfAut}}", from=2-1, to=2-2]
    \end{tikzcd}\]
\end{prop}

\begin{proof}
    We first show that the composition $\sfSph(n)\inclto \sfAdj(n)\to \sfEnd(n)$ factors through $\sfAut(n)$. This is trivial for $n=0,1$. For $n=2$, this follows from one direction of Lemma~\ref{lem:twocatsspherical}. For $n>2$, this follows by induction using that $\sfAut(-)$ is 2-Segal. So the map $\pi:\sfAdj(-)\simeq\sfBdj(-)\to\sfEnd(-)$ of Proposition~\ref{prop:simplicialbdjtoend} induces the desired map.

    By Proposition~\ref{prop:activecartesian4}, $\pi(-):\sfAdj(-)\to\sfEnd(-)$ is Cartesian over $\Delta_{\act}$. We deduce that over $a:[1]\to[2]$, the natural map
    \[\sfSph(2)\to\sfAut(2)\times_{\sfAut}\sfSph\]
    is 2-fully faithful. The other direction of Lemma~\ref{lem:twocatsspherical} proves that this map is essentially surjective, proving the desired result for $a:[1]\to[2]$. Let $b:[n-1]\to [n]$ be any active arrow. The 2-Segal property gives us a map of pullback squares of $\two$-categories as follows, where the left vertical arrows come from $b$ and the right from $a$.
    \[\begin{tikzcd}[ampersand replacement=\&]
        {\sfSph(n)}\pullback \& {\sfSph(2)} \\
        {\sfSph(n-1)} \& {\sfSph(1)}
        \arrow[from=1-1, to=1-2]
        \arrow[from=1-1, to=2-1]
        \arrow[from=1-2, to=2-2]
        \arrow[from=2-1, to=2-2]
    \end{tikzcd}
    \to
    \begin{tikzcd}[ampersand replacement=\&]
        {\sfAut(n)}\pullback \& {\sfAut(2)} \\
        {\sfAut(n-1)} \& {\sfAut(1)}
        \arrow[from=1-1, to=1-2]
        \arrow[from=1-1, to=2-1]
        \arrow[from=1-2, to=2-2]
        \arrow[from=2-1, to=2-2]
    \end{tikzcd}\]
    Moreover, this map is Cartesian over the right vertical arrow, so by a diagram chase, it is Cartesian over the left arrow.

    This proves that $\sfSph(-)\to\sfAut(-)$ is Cartesian over $\Delta_{\act}$.
\end{proof}

\subsection{Braid group action on abstract schobers}\label{subsection:mutationsph}
We now construct a $\two$-categorical braid group action on $\sfSph(n)$. We note that a $\two$-categorical braid group action has already been announced to be in the upcoming work of Abell\'an--Christ--Jasso \cite{acj}. Their approach considers all Lagrangian skeleta of $(\bbC,R)$ which includes all systems of cuts, as well as more general skeleta where multiple arcs can be incident to one singularity. In comparison, we will only consider those skeleta in $\Cuts(\bbC,R)$ (see Definition~\ref{defn:cuts}), but the novelty is that this action will be evidently compatible with abstract localized schobers.

\begin{cor}\label{cor:action2}
    There is a natural action of the braid group $\Br_n$ on $\sfSph(n)$, which respects the pullback square of Proposition~\ref{prop:activecartesian2}. Moreover, this agrees at the level of objects with the action defined in \cite{kapranovPerverseSchobers2015}*{(2.8)}.
\end{cor}

\begin{proof}
    We can pullback the action of $\Br_n$ on $\sfAut(n)$ via the pullback square of Proposition~\ref{prop:activecartesian2}, where we take $\Br_n$ to act trivially on $\sfAut$ and $\sfSph$.
    
    We compute the resulting braid group action on $\sfSph(2)$. Let $(\Phi_i,\Psi,S_i,R_i)\in\sfSph(2)$. The action of the negative generator $\sigma^{-1}\in\Br_2$ on the corresponding element of $\sfAut(2)\times_{\sfAut}\sfSph$ is as follows.
    % https://q.uiver.app/#q=WzAsMTAsWzAsMiwiXFxQaGlfMSJdLFsxLDIsIlxcUGhpXzEiXSxbMSwxLCJcXFBoaSJdLFsxLDAsIlxcUHNpIl0sWzIsMV0sWzMsMV0sWzMsMiwiXFxQaGlfMiJdLFs0LDIsIlxcUGhpXzEiXSxbNCwxLCJcXFBoaSJdLFs0LDAsIlxcUHNpIl0sWzAsMiwiSV8xIl0sWzEsMiwiSV8yIiwyXSxbMiwzLCJTIl0sWzQsNSwiXFxzaWdtYV57LTF9XFxjZG90LSIsMCx7InN0eWxlIjp7InRhaWwiOnsibmFtZSI6Im1hcHMgdG8ifX19XSxbNiw4LCJUSV8yIl0sWzcsOCwiSV8xIiwyXSxbOCw5LCJTIl1d
    \[\begin{tikzcd}[ampersand replacement=\&]
        \& \Psi \&\&\& \Psi \\
        \& \Phi \& {} \& {} \& \Phi \\
        {\Phi_1} \& {\Phi_1} \&\& {\Phi_2} \& {\Phi_1}
        \arrow["S", from=2-2, to=1-2]
        \arrow["{\sigma^{-1}\cdot-}", maps to, from=2-3, to=2-4]
        \arrow["S", from=2-5, to=1-5]
        \arrow["{I_1}", from=3-1, to=2-2]
        \arrow["{I_2}"', from=3-2, to=2-2]
        \arrow["{TI_2}", from=3-4, to=2-5]
        \arrow["{I_1}"', from=3-5, to=2-5]
    \end{tikzcd}\]
    The only non-trivial part is expressing the spherical functor $STI_2$ in terms of $S_i$ and $R_i$. Writing $U:=\Cone(SR\to\id)$, we have the identification $ST\simeq US$. We have the factorization $U\simeq U_1U_2$, where $U_i:=\Cone(S_iR_i\to\id)$, which also came up in the proof of Lemma~\ref{lem:twocatsspherical}. So
    \[STI_2\simeq USI_2 \simeq U_1U_2S_2 \simeq U_1S_2T_2.\]
    This agrees with \cite{kapranovPerverseSchobers2015}*{(2.8)} as claimed. Note that they write $T_i$ for the autoequivalence on $\Psi$ that we call $U_i$.

    The proof for general $n$ follows by the same inductive scheme as the one used to define the braid group action on $\sfAut(n)$ in Theorem~\ref{thm:inftyadmissible}.
\end{proof}

By using the same construction as Definition~\ref{defn:nonabstract}, we may define the $\two$-category $\sftwoP(\bbC,R)$ of perverse schobers.

\subsection{Fourier transforms and Stokes structures}\label{subsection:fourier}
In \cite{kapranovPerverseSchobersAlgebra2020}*{\S 4.1}, a Fourier transform for perverse schobers on $(\bbC_z,R)$ is developed (assuming that $R$ satisfies a genericity property). This is given by a map on the level of objects,
\begin{align*}
    \opname{FT}:\sftwoP(\bbC_z,R) &\isomto \sftwoP^{\opname{Stokes}}(\bbC_w,0)\\
    \calF&\mapsto\check{\calF}.
\end{align*}
The target is the collection of perverse schobers equipped with $R$-Stokes data \cite{kapranovPerverseSchobersAlgebra2020}*{\S 4.1.D}, a categorification of Stokes data which appears in the irregular Riemann-Hilbert correspondence. This Stokes data is roughly a $\zeta=\opname{Arg}(w)\in S^1$-family of SODs of $\check{\calF}\resto{\bbC^\times_w}$.

We recast this Fourier transform as an equivalence of $\two$-categories in our setting. When $R=\{0\}$, there are no Stokes structures to worry about and the Fourier transform roughly interchanges $\Phi$ and $\Psi$ (See for example \cite{gammagePerverseSchobers3d2023}*{Definition 2.3}). This identifies the projection functor $\sftwoP(\bbC_z,0)\to\ol{\sftwoP}(\bbC_z,0)$ with the restriction to the complement of $0\in\bbC_w$, as expressed by the following commuting square.
\[\begin{tikzcd}
    {\sftwoP(\bbC_z,0)}\pullback & {\sftwoP(\bbC_w,0)} \\
    {\ol{\sftwoP}(\bbC_z,0)} & {\sftwoL(\bbC_w^\times)}
    \arrow["{\opname{FT}}", from=1-1, to=1-2]
    \arrow["\pi"', from=1-1, to=2-1]
    \arrow["-|_{\bbC_w^\times}"', from=1-2, to=2-2]
    \arrow["\opname{FT}"', from=2-1, to=2-2]
\end{tikzcd}\]
Thus under Fourier transform, localized schobers on $(\bbC_z,0)$ correspond to local systems on $\bbC_w^\times$.

For general $R\subset \bbC$, we may combine this with the square from Theorem~\ref{thm:merge}.
\[\begin{tikzcd}
    {\sftwoP(\bbC_z,R)}\pullback & {\sftwoP(\bbC_z,0)}\pullback & {\sftwoP(\bbC_w,0)} \\
    {\ol{\sftwoP}(\bbC_z,R)} & {\ol{\sftwoP}(\bbC_z,0)} & {\sftwoL(\bbC_w^\times)}
    \arrow["{\FS_{\sftwoP}}", from=1-1, to=1-2]
    \arrow["\pi"', from=1-1, to=2-1]
    \arrow["{\opname{FT}}", from=1-2, to=1-3]
    \arrow["\pi"', from=1-2, to=2-2]
    \arrow["-|_{\bbC_w^\times}"', from=1-3, to=2-3]
    \arrow["{\FS_{\ol{\sftwoP}}}"', from=2-1, to=2-2]
    \arrow["\opname{FT}"', from=2-2, to=2-3]
\end{tikzcd}\]
The outer pullback square gives us an equivalence,
\[\sftwoP(\bbC_z,R)\isomto \sftwoP(\bbC_w,0)\times_{\sftwoL(\bbC_w^\times)}\ol{\sftwoP}(\bbC_z,R).\]
This fiber product can be interpreted as equipping the nearby cycles part of an object of $\sftwoP(\bbC_w,0)$ with an $R$-Stokes structure. Thus under Fourier transform, we may understand a localized schobers on $(\bbC,R)$ as none other than a local system on $\bbC_w-0$ equipped with an $R$-Stokes structure.